\section{Test Trained Agent on Different Time Periods}

The data is divided into train sets and test sets according to the trading time. After the
agents are trained over a specified timeframe, they are then evaluated on the test set
consisting of a different timeframe.

\begin{figure}[H]
  \centering
  \includegraphics[width=0.78\textwidth]{z129_1.png}
  \caption{Train/test split -- cumulative return ratio (split 1).}
\end{figure}

\begin{figure}[H]
  \centering
  \includegraphics[width=0.78\textwidth]{z130_1.png}
  \caption{Train/test split -- cumulative return ratio (split 2).}
\end{figure}

\begin{figure}[H]
  \centering
  \includegraphics[width=0.78\textwidth]{z132_1.png}
  \caption{Train/test split -- cumulative return ratio (split 4).}
\end{figure}

\paragraph{Comments.}
Results show that even when the trained agent shows promising results in terms of profits,
if applied on the test set (except for one trained agent), positive rewards turn into
negative rewards. Even if they show consistent behavior, the outcome isn't.

As a central conclusion, two things can be stated:
\begin{itemize}
  \item If the agent is only trained over 4 years, no reliable predictions can be made
        about future behavior.
  \item The algorithm used here offers no possibility of predictability for future
        behavior in the field of stock trading and, on the contrary, performs worse than
        simple buy and hold.
\end{itemize}
